\section{Análisis del Prompt del Paso 1: generación de la solución inicial}

\subsection{Estructura del Prompt}

El System Prompt del Paso 1 es un texto en Markdown estructurado estándar, dividido en tres partes:

\begin{enumerate}
    \item \textbf{Instrucciones centrales}:
    \begin{itemize}
        \item El rigor es fundamental (la mayoría de los problemas de la IMO son de demostración, y el criterio final es el rigor lógico)
        \item Sé honesto respecto al resultado: no lo fuerces (para evitar el Reward Hacking)
        \item Usa LaTeX en todo el texto
    \end{itemize}
    \item \textbf{Formato de salida (Adaptive Prompt)}:
    \begin{itemize}
        \item Summary: incluye el Verdict (si se encontró una solución completa) y un resumen del método
        \item Detailed Solution: la solución detallada
    \end{itemize}
    \item \textbf{Instrucciones de Self-Verification}
\end{enumerate}

\subsection{Intención de diseño del formato de salida}

\begin{importantbox}{El Verdict se usa para decidir la ramificación}
El Verdict dentro del Summary no solo presenta información, sino que además es \textbf{el criterio para decidir la ramificación} del Workflow posterior: si el Verdict es false (no se encontró una solución completa), se omite la verificación posterior y se vuelve a hacer sample directamente. La Detailed Solution, en cambio, se envía a la etapa de Verification.
\end{importantbox}

El código de este texto usa \textbf{expresiones regulares} para extraer estas salidas estructuradas (en lugar de una API de salida estructurada), porque la capacidad de seguimiento de instrucciones de Gemini 2.5 Pro es lo bastante sólida.

\subsection{Paso 2: Self-Improvement}

El Paso 2 \textbf{comparte el contexto} con el Paso 1. Su Prompt es sencillo: «tienes una oportunidad de mejorar tu resultado: haz un review cuidadoso de tu solution, corrige errores y completa los justification gap».

La razón para introducir el Paso 2 es que, en el Paso 1, el Thinking Budget (32768 tokens) puede agotarse por completo, lo que obliga al modelo a terminar de pensar de forma apresurada. El Paso 2 le da al modelo la oportunidad de seguir pensando.

\subsection{Resumen de la sección}

Lo esencial del diseño del Prompt es que la salida estructurada sienta las bases para decidir la ramificación y el flujo de datos del Workflow posterior.
